\section{III. RESULT}

\subsection{Paragraf 1: Arsitektur Protokol}
\begin{frame}{Paragraf 1}
    \begin{itemize}
        \item \textbf{1. pembagian fase:} Protokol dibagi menjadi dua fase utama untuk solusi spektral komprehensif.
        \item \textbf{2. estimasi eigenvector:} Prioritas pertama adalah estimasi $\ket{u_{max}}$ berdasarkan $\lambda_{max}$.
        \item \textbf{3. alur evolusioner:} Inisialisasi acak $\to$ Sirkuit Fig. 2 $\to$ Proyeksi $n$-bit $\to$ Pendekatan bertahap ke vektor eksak.
        \item \textbf{4. input QPE:} \textit{Eigenvector} teroptimasi digunakan sebagai input untuk ekstraksi \textit{eigenvalue} presisi tinggi.
    \end{itemize}
\end{frame}

\begin{frame}{Pertanyaan yang dijawab pada Paragraf 1}
    \centering
    \small
    \begin{tabular}{lp{4.5cm}p{6cm}}
        \toprule
        Kalimat ke & Pertanyaan & Jawaban \\
        \midrule
        1 & Struktur protokol? & Dibagi menjadi estimasi \textit{eigenvector} dan perbaikan \textit{eigenvalue}. \\
        2 & Fokus utama bagian 1? & Mengestimasi $\ket{u_{max}}$ yang berkorespondensi dengan $\lambda_{max}$. \\
        4 & Fungsi \textit{eigenvector}? & Sebagai \textit{initial state} QPE untuk presisi \textit{eigenvalue} lebih tinggi. \\
        \bottomrule
    \end{tabular}
\end{frame}

\subsection{Paragraf 2: Prosedur Estimasi Eigenvector}
\begin{frame}{Paragraf 2}
    \begin{itemize}
        \item \textbf{1. titik awal:} Inisialisasi menggunakan \textit{random state} kuantum $\ket{b_0}$.
        \item \textbf{2. konfigurasi register:} \textit{Initial state} sistem didefinisikan sebagai $\ket{0} \otimes \ket{0} \otimes \ket{b_0}$.
        \item \textbf{3. strategi feedback:} Hasil estimasi $\ket{b_1}$ digunakan kembali sebagai state awal iterasi berikutnya.
        \item \textbf{5. kriteria berhenti:} Iterasi berlanjut hingga tercapai stabilitas $\ket{b_{k-1}} \approx \ket{b_k}$.
        \item \textbf{6. klaim konvergensi:} State $\ket{b_k}$ diklaim sebagai aproksimasi terbaik untuk $\ket{u_{max}}$.
    \end{itemize}
\end{frame}

\begin{frame}{Pertanyaan yang dijawab pada Paragraf 2}
    \centering
    \small
    \begin{tabular}{lp{4.5cm}p{6cm}}
        \toprule
        Kalimat ke & Pertanyaan & Jawaban \\
        \midrule
        1 & Titik awal estimasi? & Dimulai dengan \textit{random state} $\ket{b_0}$. \\
        3 & Proses naikkan fidelitas? & Hasil $\ket{b_1}$ digunakan kembali sebagai state inisial iterasi berikutnya. \\
        5 & Kapan iterasi selesai? & Dilakukan $k$ kali hingga tercapai stabilitas $\ket{b_{k-1}} \approx \ket{b_k}$. \\
        \bottomrule
    \end{tabular}
\end{frame}

\subsection{Paragraf 3: Penentuan Eigenvalue Numerik}
\begin{frame}{Paragraf 3}
    \begin{itemize}
        \item \textbf{1. target kedua:} Estimasi nilai numerik \textit{eigenvalue} $\lambda_{max}$ setelah basis ditemukan.
        \item \textbf{2. aplikasi QPE:} Menggunakan algoritma QPE standar untuk mendapatkan presisi $n$-bit.
        \item \textbf{3. limit presisi:} Akurasi dibatasi secara pragmatis oleh kapasitas fisik prosesor kuantum.
        \item \textbf{5. urutan penyelesaian:} Strategi \textit{bottom-up} (kasus $2 \times 2$ sebelum ekspansi $4 \times 4$).
        \item \textbf{7. konsistensi hasil:} Laporan hasil yang akurat antara simulator Qiskit dan hardware riil.
    \end{itemize}
\end{frame}

\begin{frame}{Pertanyaan yang dijawab pada Paragraf 3}
    \centering
    \small
    \begin{tabular}{lp{4.5cm}p{6cm}}
        \toprule
        Kalimat ke & Pertanyaan & Jawaban \\
        \midrule
        1 & Target pasca estimasi? & Estimasi nilai numerik \textit{eigenvalue} $\lambda_{max}$ secara akurat. \\
        3 & Faktor pembatas presisi? & Ukuran prosesor membatasi jumlah bit presisi yang dicapai. \\
        5 & Urutan penyelesaian? & Sub-matriks $2 \times 2$ lalu beralih ke ekspansi $4 \times 4$. \\
        \bottomrule
    \end{tabular}
\end{frame}

\subsection{Paragraf 4: Normalisasi dan Spektral PCA}
\begin{frame}{Paragraf 4}
    \begin{itemize}
        \item \textbf{1. formalisme density matrix:} Normalisasi matriks kovarians terhadap \textit{trace} untuk representasi kuantum.
        \item \textbf{2. dominasi spektral:} Identifikasi $\lambda_{max} \gg \lambda_2$ sebagai fitur intrinsik matriks korelasi finansial.
        \item \textbf{3. uniter evolusi:} Definisi operator uniter berdasarkan dekomposisi spektral matriks ternormalisasi.
    \end{itemize}
\end{frame}

\begin{frame}{Pertanyaan yang dijawab pada Paragraf 4}
    \centering
    \small
    \begin{tabular}{lp{4.5cm}p{6cm}}
        \toprule
        Kalimat ke & Pertanyaan & Jawaban \\
        \midrule
        1 & Persiapan matriks? & Normalisasi terhadap \textit{trace} agar sesuai \textit{density matrix}. \\
        2 & Mengapa PCA relevan? & Karena \textit{eigenvalue} utama jauh lebih besar dari yang kedua ($\lambda_{max} \gg \lambda_2$). \\
        \bottomrule
    \end{tabular}
\end{frame}

\subsection{Paragraf 5: Konvergensi dan Pengukuran Basis}
\begin{frame}{Paragraf 5}
    \begin{itemize}
        \item \textbf{1. alokasi qubit:} 2 qubit untuk estimasi \textit{eigenvalue} (2-bit) dan 1 qubit untuk \textit{eigenvector}.
        \item \textbf{3. stabilitas sistem:} Vektor hasil pengukuran mencapai titik stabil setelah iterasi ke-4.
        \item \textbf{4. rotasi basis:} Penggunaan basis x, y, dan $r$ untuk mengekstraksi informasi fasa kompleks.
        \item \textbf{7. metrik galat:} Definisi $\delta$ sebagai batas atas seluruh kesalahan eksperimental sistematis.
        \item \textbf{12. observasi:} Algoritma menunjukkan tanda konvergensi signifikan sejak iterasi pertama.
    \end{itemize}
\end{frame}

\begin{frame}{Pertanyaan yang dijawab pada Paragraf 5}
    \centering
    \small
    \begin{tabular}{lp{4.5cm}p{6cm}}
        \toprule
        Kalimat ke & Pertanyaan & Jawaban \\
        \midrule
        1 & Alokasi qubit? & 3 qubit: 2 untuk $\lambda$ (2-bit) dan 1 untuk $\ket{u}$. \\
        4 & Mengapa rotasi basis? & Mendapatkan informasi fasa kompleks yang tidak terbaca pada basis standar. \\
        7 & Fungsi parameter $\delta$? & Metrik ketidakpastian yang mencakup berbagai sumber galat sistematis. \\
        12 & Temuan konvergensi? & Algoritma sudah menunjukkan tanda konvergensi signifikan sejak iterasi pertama. \\
        \bottomrule
    \end{tabular}
\end{frame}

\subsection{Paragraf 6: Strategi Pembagian Tahapan (Stages)}
\begin{frame}{Paragraf 6}
    \begin{itemize}
        \item \textbf{1. estimasi awal:} Hasil diperoleh melalui proyeksi pada \textit{subspace} dengan resolusi 2-bit.
        \item \textbf{3. rasionalitas tahap:} Pembagian protokol untuk optimalisasi sumber daya dan minimisasi galat.
        \item \textbf{4. alasan heuristik:} Ketidaktahuan nilai eigen secara \textit{a priori} membutuhkan pencarian awal.
        \item \textbf{5. akumulasi galat:} Observasi empiris menunjukkan peningkatan galat jika dijalankan satu tahap.
    \end{itemize}
\end{frame}

\begin{frame}{Pertanyaan yang dijawab pada Paragraf 6}
    \centering
    \small
    \begin{tabular}{lp{4.5cm}p{6cm}}
        \toprule
        Kalimat ke & Pertanyaan & Jawaban \\
        \midrule
        3 & Mengapa bagi 2 tahap? & Berdasarkan alasan strategis: ketidaktahuan awal dan efisiensi hardware. \\
        5 & Alasan observasi? & Observasi akumulasi kesalahan gerbang yang meningkat jika dijalankan sekaligus. \\
        \bottomrule
    \end{tabular}
\end{frame}

\subsection{Paragraf 7: Resolusi dan Dekoherensi}
\begin{frame}{Paragraf 7}
    \begin{itemize}
        \item \textbf{1. peningkatan resolusi:} Menggunakan 3 qubit untuk mengevaluasi batas kemampuan hardware.
        \item \textbf{2. fenomena dekoherensi:} Laporan dekoherensi signifikan akibat pertambahan kedalaman sirkuit.
        \item \textbf{3. fungsi simulator:} Simulator Qiskit digunakan sebagai standar pembanding fidelitas ideal.
        \item \textbf{5. prospek masa depan:} Perbaikan hasil bergantung pada fidelitas gerbang dan kualitas chip.
    \end{itemize}
\end{frame}

\begin{frame}{Pertanyaan yang dijawab pada Paragraf 7}
    \centering
    \small
    \begin{tabular}{lp{4.5cm}p{6cm}}
        \toprule
        Kalimat ke & Pertanyaan & Jawaban \\
        \midrule
        2 & Dampak kedalaman? & Munculnya dekoherensi signifikan (Gbr. 3). \\
        3 & Kegunaan simulator? & Standar pembanding untuk menyoroti tingkat derau pada hardware riil. \\
        5 & Prospek perbaikan? & Bergantung peningkatan fidelitas gerbang dan kualitas chip kuantum. \\
        \bottomrule
    \end{tabular}
\end{frame}

\subsection{Paragraf 8: Ekspansi Kasus Matriks 4x4}
\begin{frame}{Paragraf 8}
    \begin{itemize}
        \item \textbf{1. skalabilitas:} Uji algoritma pada matriks $4 \times 4$ untuk skenario finansial luas.
        \item \textbf{5. reduksi kedalaman:} Dekomposisi gerbang menjadi interaksi 2-qubit untuk menekan dekoherensi.
        \item \textbf{6. desain sirkuit:} Referensi Gambar 4 untuk representasi gerbang teroptimasi kasus $4 \times 4$.
        \item \textbf{7. inisialisasi state:} Penggunaan superposisi merata $\ket{b_0}$ sebagai titik awal pencarian.
    \end{itemize}
\end{frame}

\begin{frame}{Pertanyaan yang dijawab pada Paragraf 8}
    \centering
    \small
    \begin{tabular}{lp{4.5cm}p{6cm}}
        \toprule
        Kalimat ke & Pertanyaan & Jawaban \\
        \midrule
        1 & Mengapa ke $4 \times 4$? & Menguji skalabilitas algoritma pada skenario finansial lebih kompleks. \\
        5 & Cara tekan dekoherensi? & Dekomposisi sirkuit menjadi interaksi 2-qubit (\textit{reduced depth}). \\
        \bottomrule
    \end{tabular}
\end{frame}

\subsection{Paragraf 9: Evaluasi Akumulasi Galat 4x4}
\begin{frame}{Paragraf 9}
    \begin{itemize}
        \item \textbf{1. objektivitas data:} Pengujian pada berbagai basis untuk menghitung rata-rata kesalahan.
        \item \textbf{3. batas toleransi:} Pengakuan bahwa akumulasi galat hardware riil melampaui batas untuk hasil konklusif.
        \item \textbf{analisis induktif:} Hasil mentah menunjukkan perlunya metodologi mitigasi tingkat lanjut.
    \end{itemize}
\end{frame}

\begin{frame}{Pertanyaan yang dijawab pada Paragraf 9}
    \centering
    \small
    \begin{tabular}{lp{4.5cm}p{6cm}}
        \toprule
        Kalimat ke & Pertanyaan & Jawaban \\
        \midrule
        1 & Jaminan objektivitas? & Menguji hasil pada berbagai basis berbeda dan menghitung rata-rata galat. \\
        3 & Mengapa tidak konklusif? & Akumulasi galat hardware riil melampaui ambang batas toleransi. \\
        \bottomrule
    \end{tabular}
\end{frame}

\subsection{Paragraf 10: Metodologi Mitigasi Kesalahan}
\begin{frame}{Paragraf 10}
    \begin{itemize}
        \item \textbf{1. referensi mitigasi:} Penerapan teknik perbaikan data berdasarkan literatur terbaru (Ref [41]).
        \item \textbf{3. ekstrapolasi matematis:} Penggunaan \textit{Richardson's extrapolation} untuk eliminasi gangguan derau.
        \item \textbf{4. readout mitigation:} Lapisan pelindung tambahan terhadap galat pada tahap pengukuran.
    \end{itemize}
\end{frame}

\begin{frame}{Pertanyaan yang dijawab pada Paragraf 10}
    \centering
    \small
    \begin{tabular}{lp{4.5cm}p{6cm}}
        \toprule
        Kalimat ke & Pertanyaan & Jawaban \\
        \midrule
        1 & Langkah perbaikan data? & Menerapkan teknik mitigasi galat dari literatur (Ref [41]). \\
        3 & Teknik ekstrapolasi? & \textit{Richardson's extrapolation} untuk eliminasi derau secara matematis. \\
        4 & Lapisan mitigasi ekstra? & \textit{Readout error mitigation} untuk mereduksi galat spesifik pengukuran. \\
        \bottomrule
    \end{tabular}
\end{frame}
